% 《大学物理学学习指导》第5版，自测题11填空11.2.1—11.2.9。
% 源书印刷231—232页，PDF物理239—240页。
\subsection{自测题11：光谱、粒子波长和量子态填空}

\sourcetest{11.2.1}{231}{239}
\problemtext X射线散射后存在两种波长分量，
哪一种称为康普顿散射分量？
\workings
散射光包括与入射波长\textbf{相同}的弹性散射成分，
也包括波长比入射光\textbf{更长}的非弹性散射成分。
后者来自电子反冲的能量转移，称
这就是所填的\textbf{康普顿散射分量}。
\crosscheck 无自由电子反冲时的相干弹性散射不能
因为同在散射光束中就一并称为“康普顿位移”。

\sourcetest{11.2.2}{232}{240}
\problemtext 氢原子巴耳末线系最长波长与最短波长的比值。
\workings
终态$n=2$。
最长线为$3\to2$，
最短波长为系列极限$\infty\to2$。
用里德伯关系
\[
\boxed{\frac{\lambda_{\max}}{\lambda_{\min}}
=\frac{1/4}{1/4-1/9}
=\frac95=1.80}.
\]
\crosscheck 最短波长是\emph{巴耳末极限}，
不是从有限$n=4$的次长线误推得到的27/20。

\sourcetest{11.2.3}{232}{240}
\problemtext 氢原子电子位于$n=3$时，
完全电离到自由电子所需能量是多少eV？
\workings
玻尔氢原子能级$E_n=-13.6/n^2$eV，
电离极限$E_\infty=0$，
\[
\boxed{E_{\rm ion}=0-E_3=\frac{13.6}{9}
\simeq1.51\,\mathrm{eV}}.
\]
\crosscheck 不可按基态电离能13.6eV直接填写，
因为原电子已经在$n=3$激发态。

\sourcetest{11.2.4}{232}{240}
\problemtext 中子德布罗意波长$0.2$nm，
取$h=6.63\times10^{-34}$J\,s、
$m=1.67\times10^{-27}$kg，
求其非相对论动能。
\workings
$p=h/\lambda$，$K=p^2/(2m)$，故
\[
\boxed{K=\frac{h^2}{2m\lambda^2}
=\frac{(6.63\times10^{-34})^2}
{2(1.67\times10^{-27})(0.2\times10^{-9})^2}
\simeq3.29\times10^{-21}\,\mathrm J}.
\]
折合约$0.0205$eV。
\crosscheck 中子的静止能量约939MeV，
本题动能极低，非相对论动能公式充分适用。

\sourcetest{11.2.5}{232}{240}
\problemtext 电子局域于长度$\Delta x=0.05$nm范围，
按原书指定$\Delta x\Delta p_x\ge h$估算
$x$方向动量不确定度。
\workings
原书此题明确采用$h$而非$\hbar/2$，
\[
\boxed{\Delta p_x\gtrsim
\frac{6.63\times10^{-34}}
{0.05\times10^{-9}}
\simeq1.33\times10^{-23}\,\mathrm{kg\,m/s}}.
\]
\crosscheck 若用标准差形式
$\Delta x\Delta p_x\ge\hbar/2$，
下界应再除以$4\pi$；
两种约定的结果不应混称为同一个精确界限。

\sourcetest{11.2.6}{232}{240}
\problemtext 补全三组$n=2$单电子四量子数中的合法缺项：
(1) $m_l=-1,m_s=-1/2$；
(2) $l=0,m_s=+1/2$；
(3) $l=1,m_l=0$。
\workings
量子数允许范围
$l=0,\ldots,n-1$，
$m_l=-l,\ldots,l$，
$m_s=\pm1/2$。
因此
\[
\boxed{\text{(1)}\ l=1},\qquad
\boxed{\text{(2)}\ m_l=0},\qquad
\boxed{\text{(3)}\ m_s=\pm\frac12}.
\]
\crosscheck $m_l=-1$要求$l\ge1$，
与$n=2$一起唯一确定$l=1$。
原题(3)只给出一个没有其他限制的电子，
允许两种自旋投影而非只能取某一符号。

\sourcetest{11.2.7}{232}{240}
\problemtext 主量子数$n=4$时角量子数$l$
及磁量子数$m_l$可能取值。
\workings
$n=4$时$l=0,1,2,3$。
对每个$l$都有$m_l=-l,\ldots,l$，
综合所有允许$l$，
\[
\boxed{l=0,1,2,3},\qquad
\boxed{m_l=-3,-2,-1,0,1,2,3}.
\]
\crosscheck 虽然总集合含7个$m_l$值，
但并不是每个$l$均可同时取这7个值；
必须满足$|m_l|\le l$。

\sourcetest{11.2.8}{232}{240}
\problemtext 基态氦原子的两名电子用四量子数
$(n,l,m_l,m_s)$分别如何表示？
\workings
基态氦的两个电子占据$1s$轨道；
泡利不相容原理要求自旋投影相反，
所以两组为
\[
\boxed{(1,0,0,+1/2)},\qquad
\boxed{(1,0,0,-1/2)}.
\]
\crosscheck 相同的前三量子数是允许的；
不允许两个电子的\emph{四个}量子数同时完全相同。

\sourcetest{11.2.9}{232}{240}
\problemtext 原子序数$Z=6$的碳及$Z=14$的硅，
基态电子组态分别是什么？
\workings
在原书采用的独立电子填充近似下，
按照能量顺序和泡利不相容原理，
\[
\boxed{\mathrm C:\ 1s^2\,2s^2\,2p^2},
\]
\[
\boxed{\mathrm{Si}:\ 1s^2\,2s^2\,2p^6\,3s^2\,3p^2}.
\]
\crosscheck 电子占据数分别为$2+2+2=6$，
及$2+2+6+2+2=14$；
开放$p$亚层的具体自旋排列还受洪特规则影响。
